% H08-S001 | sumber authority/habring/source-v1/stochastic.tex baris 1--34
% segment-id: d90.hab.v1.ch08.seg0001
\chapter{Penurunan Gradien Stokastik}
Andaikan kita ingin menyelesaikan
\begin{equation}
    \min_x f(x),
\end{equation}
tetapi hanya mempunyai akses ke estimasi acak yang tak bias bagi suatu subgradien $f$. Artinya, kita mengandaikan dapat mengevaluasi $G(x,z)$ yang memenuhi
\begin{equation}\label{stochastic:eq:gradient}
    \E[G(x,Z)] \in \partial f(x),
\end{equation}
dengan $Z$ suatu peubah acak.
Ternyata banyak bukti konvergensi dapat dialihkan ke latar stokastik ini dengan hampir tanpa upaya tambahan. Sebelum menunjukkannya, mari kita berikan motivasi untuk latar ini.

Dalam pembelajaran mesin modern, banyak masalah pembelajaran dapat dirumuskan sebagai masalah jumlah hingga. Agar pengambilan sampel seragam menghasilkan estimator gradien yang tak bias, kita memakai normalisasi rerata
\begin{equation}
    \min_x f(x)\coloneqq \frac{1}{N}\sum_{i=1}^N f_i(x),
\end{equation}
dengan $N$ biasanya merupakan banyaknya sampel pelatihan. Sebagai contoh, andaikan kita ingin mempelajari pemetaan berparameter $f_\theta$ yang seharusnya menghampiri relasi $x\mapsto y$ berdasarkan data pelatihan $(x_i,y_i)_{i=1}^N$. Masalah pelatihannya dapat berupa
\begin{equation}
    \min_\theta \frac{1}{N}\sum_{i=1}^N \|y_i-f_\theta(x_i)\|^2.
\end{equation}
Karena ukuran data pelatihan yang dijumpai, sering kali gradien penuh $f$ tidak dapat dihitung secara langsung akibat keterbatasan memori.
Jika semua $f_i$ terdiferensialkan, salah satu jalan keluarnya adalah pembaruan stokastik
\begin{equation}
    \begin{cases}
        \text{Pilih $i\in\Ic$ secara seragam}\\
        x_{k+1} = x_k - \tau \nabla f_i(x_k)
    \end{cases}
\end{equation}
dengan $\Ic=\{1,\dots,N\}$.
Dengan mendefinisikan $(Z_k)_k$ sebagai barisan peubah acak iid yang berdistribusi seragam pada $\Ic$, pembaruan di atas dapat ditulis sebagai
\begin{equation}
    x_{k+1} = x_k - \tau G(x_k,Z_k)
\end{equation}
dengan $G(x,z) \coloneqq \nabla f_z(x)$. Memang,
$\E[G(x,Z)]=N^{-1}\sum_{i=1}^N\nabla f_i(x)=\nabla f(x)$, sehingga $G$ memenuhi~\eqref{stochastic:eq:gradient}.

% H08-S002 | sumber authority/habring/source-v1/stochastic.tex baris 36--50
% segment-id: d90.hab.v1.ch08.seg0002
\begin{theorem}
    Misalkan $C\subseteq\R^d$ tak kosong, tertutup, dan konveks, serta misalkan $f:\R^d\rightarrow\R$ konveks dan $L$-Lipschitz. Andaikan masalah $\min_{x\in C}f(x)$ mempunyai solusi, pilih $x^*\in\arg\min_{x\in C}f(x)$, dan ambil $x_0\in C$. Untuk $k\geq0$, misalkan
    $G_k\coloneqq G(x_k,Z_k)$ dan jalankan iterasi terproyeksi
    $x_{k+1}=\proj_C(x_k-\tau_kG_k)$ dengan $\tau_k>0$.

    Tuliskan $\mathcal F_k=\sigma(x_0,Z_0,\dots,Z_{k-1})$, sehingga $x_k$ terukur terhadap $\mathcal F_k$, dan definisikan $\bar g_k\coloneqq\E[G_k\mid\mathcal F_k]$. Andaikan, hampir pasti untuk setiap $k$,
    \begin{equation}
        \begin{cases}
            \bar g_k \in \partial f(x_k),\\
            \E[\|G_k-\bar g_k\|^2\mid\mathcal F_k]\leq \sigma^2<\infty.
        \end{cases}
    \end{equation}
    Andaikan pula ukuran langkah memenuhi
    \[
        \sum_{k=0}^{K-1}\tau_k\longrightarrow\infty,
        \qquad
        \frac{\sum_{k=0}^{K-1}\tau_k^2}{\sum_{k=0}^{K-1}\tau_k}\longrightarrow0
        \quad\text{ketika }K\longrightarrow\infty.
    \]
    Untuk $K\geq1$, definisikan
    $f_{\mathrm{best}}^K\coloneqq\min_{0\leq k\leq K-1}f(x_k)$.
    Maka $\E[f_{\mathrm{best}}^K-f(x^*)]\rightarrow0$ ketika $K\rightarrow\infty$.
\end{theorem}

% H08-S003 | sumber authority/habring/source-v1/stochastic.tex baris 51--107
% segment-id: d90.hab.v1.ch08.seg0003
\begin{proof}
    Karena $x^*\in C$, kita mempunyai $\proj_C(x^*)=x^*$. Ketakmembangan proyeksi metrik memberikan
    \begin{equation}
        \begin{aligned}
            \|x_{k+1}-x^*\|^2
            &= \|\proj_C(x_{k}-\tau_k G_k)-\proj_C(x^*)\|^2\\
            &\leq \|x_{k}-\tau_k G_k-x^*\|^2\\
            &= \|x_{k}-x^*\|^2 - 2\tau_k\inner{x_k-x^*}{G_k} + \tau_k^2\|G_k\|^2.
        \end{aligned}
    \end{equation}
    Perhatikan bahwa sifat terukur $x_k$ dan ketakbiasan bersyarat memberi
    \begin{equation}
        \begin{aligned}
            \E[\inner{x_k-x^*}{G_k}]
            &= \E\!\left[\E[\inner{x_k-x^*}{G_k}\mid\mathcal F_k]\right]\\
            &= \E[\inner{x_k-x^*}{\bar g_k}]\\
            &\geq \E[f(x_k)-f(x^*)]\\
            &= \E[f(x_k)]-f(x^*),
        \end{aligned}
    \end{equation}
    dengan ketaksamaan subgradien pada baris ketiga. Karena fungsi konveks $f$ bersifat $L$-Lipschitz pada $\R^d$, setiap subgradiennya mempunyai norma paling besar $L$. Identitas momen kedua bersyarat kemudian menghasilkan
    \begin{equation}
        \E[\|G_k\|^2\mid\mathcal F_k]
        = \E[\|G_k-\bar g_k\|^2\mid\mathcal F_k]+\|\bar g_k\|^2
        \leq \sigma^2+L^2\coloneqq M^2.
    \end{equation}
    Dengan mengambil ekspektasi pada ketaksamaan pertama, diperoleh
    \begin{equation}
        \begin{aligned}
            \E[\|x_{k+1}-x^*\|^2]
            \leq \E[\|x_k-x^*\|^2]
            -2\tau_k\big(\E[f(x_k)]-f(x^*)\big)
            +\tau_k^2M^2.
        \end{aligned}
    \end{equation}
    Selebihnya mengikuti argumen pada kasus deterministik. Menjumlahkan dari $k=0$ sampai $K-1$, memakai ketaknegatifan suku jarak terakhir, lalu membagi dengan $2$ memberikan
    \begin{equation}
        \begin{aligned}
            \sum_{k=0}^{K-1}\tau_k\big(\E[f(x_k)]-f(x^*)\big)
            \leq \frac{1}{2}\|x_0-x^*\|^2
            +\frac{M^2}{2}\sum_{k=0}^{K-1}\tau_k^2.
        \end{aligned}
    \end{equation}
    Definisikan
    \begin{equation}
        S_K\coloneqq\sum_{k=0}^{K-1}\tau_k,
        \qquad
        Q_K\coloneqq\sum_{k=0}^{K-1}\tau_k^2.
    \end{equation}
    Karena $\tau_k>0$, minimum terbaik tidak melebihi rerata berbobot. Oleh karena itu,
    \begin{equation}
        \begin{aligned}
            \E[f_{\mathrm{best}}^K-f(x^*)]
            &\leq \frac{\sum_{k=0}^{K-1}\tau_k\E[f(x_k)-f(x^*)]}{S_K}\\
            &\leq \frac{\|x_0-x^*\|^2}{2S_K}
            +\frac{M^2Q_K}{2S_K}.
        \end{aligned}
    \end{equation}
    Berdasarkan kedua syarat ukuran langkah,
    \begin{equation}
        0\leq\E[f_{\mathrm{best}}^K-f(x^*)]
        \leq\frac{\|x_0-x^*\|^2}{2S_K}
        +\frac{M^2Q_K}{2S_K}
        \longrightarrow0
        \qquad(K\rightarrow\infty).
    \end{equation}
\end{proof}
